\section{智能不是一条语言分数线}

访谈结尾，谢赛宁讨论 AGI 是伪命题、动物智能和人类自大。他引用动物认知相关书籍，谈到不同生物有不同感知和行动方式：黑猩猩的推理、鸟类缓存食物、鲸鱼交流、狗的嗅觉、蝙蝠的听觉。这里的技术含义是：智能不是在一条语言考试分数线上排序，而是和身体、环境、目标、感知通道和生存任务绑定。

\begin{figure}[H]
\centering
\includegraphics[width=0.94\textwidth]{intelligence-spectrum.png}
\caption{智能不是一条语言分数线：从文本/代码到视觉、行动、身体环境和社会目标。}
\end{figure}

\begin{knowledgebox}{读图：为什么反人类中心视角重要}
图中把智能拆成文本/代码、视觉/空间、行动/控制、身体/环境和社会/目标几个维度。LLM 在文本和代码上表现强，不代表它自动拥有身体、环境和行动层面的能力。世界模型的意义，正是把智能重新放回感知、预测、行动和反馈之中。
\end{knowledgebox}

\subsection{松鼠智能与孩子家务}

访谈中引用 Rich Sutton 的一个视角：与其惊叹模型会写代码或拿数学竞赛奖牌，不如思考造出一只能够在真实世界生存的松鼠有多难。这个说法看似夸张，但它把问题从“人类考试能力”拉回“生命体在世界中行动”的基础能力。

谢赛宁进一步说，几岁孩子或十二岁孩子可以完成很多家务，而今天的机器人还无法可靠承担。这不是贬低 LLM，而是在提醒：真实世界里的感知、动作、恢复、目标和常识，仍是未解决的大问题。

\begin{warningbox}{AGI 口号容易遮蔽具体能力}
如果不定义环境、身体、目标和评测，AGI 很容易变成空泛标签。更可操作的问题是：系统能否在某个真实环境中可靠感知、预测、行动、犯错后恢复，并逐渐迁移到新任务？
\end{warningbox}

\subsection{本章小结}

智能不是单一尺度。语言能力很重要，但真实世界智能还包括视觉、身体、行动、反馈和社会目标。世界模型路线试图补的，正是这些被语言分数遮蔽的维度。

